\chapter{Conclusioni}
\label{cap:td-conclusioni}

\section{Valutazione complessiva}

I criteri di uscita definiti in \cref{cap:td-strategia} sono soddisfatti:
tutti i casi di test sono superati, la copertura supera le soglie
configurate e non restano difetti aperti. Il sistema può essere considerato
verificato rispetto ai requisiti del RAD, entro i limiti dichiarati nella
sezione seguente.

Sotto il profilo del metodo, la verifica ha prodotto due risultati che vale
la pena isolare.

Il primo riguarda i requisiti non funzionali. Averli formulati nel RAD con un
criterio di accettazione osservabile — invece che come affermazioni generiche
sulla qualità — li ha resi traducibili in asserzioni. Sette difetti di
sicurezza su quindici sono stati individuati proprio così: erano casi in cui
il sistema faceva correttamente ciò che gli veniva chiesto, ma da chiunque lo
chiedesse, e nessun test funzionale li avrebbe rivelati.

Il secondo riguarda la progettazione. La possibilità di verificare ogni
servizio in isolamento non è una proprietà della suite ma del sistema: deriva
dalla decisione, presa nel SDD, di passare le dipendenze dai costruttori.
Senza di essa la stessa verifica richiederebbe di aggirare l'incapsulamento
per sostituire un collaboratore. La difficoltà a scrivere un test si è
rivelata, in questo progetto, un indicatore attendibile di un problema di
progettazione, e in più di un caso ha portato a rivedere l'interfaccia di una
classe prima ancora di verificarla.

\section{Limiti della verifica}
\label{sec:td-limiti}

\begin{longtable}{@{}L{4.4cm} L{5.0cm} L{5.0cm}@{}}
\toprule
\textbf{Limite} & \textbf{Rischio residuo} & \textbf{Mitigazione} \\
\midrule
\endfirsthead
\toprule
\textbf{Limite} & \textbf{Rischio residuo} & \textbf{Mitigazione} \\
\midrule
\endhead
\bottomrule
\endlastfoot

Il livello di persistenza è sostituito da mock: le interrogazioni SQL non
sono eseguite dalla verifica automatica. &
Un errore in una query — nome di colonna, ordine dei parametri — non viene
intercettato prima dell'esecuzione. È accaduto: \id{DF\_12} e \id{DF\_13}
sono difetti di questa natura. &
Prove di sistema \id{TC\_SY\_01}--\id{TC\_SY\_09} su un database reale, da
eseguire prima di ogni rilascio. Il livello è sottile e privo di logica
condizionale, il che riduce la superficie esposta. \\

I trigger che mantengono la valutazione media vivono nel database. &
Una modifica ai trigger non è rilevata da alcun test automatico. &
\id{TC\_SY\_06} verifica il ricalcolo su dati reali. \\

Non esistono test automatici sull'interfaccia utente. &
Un difetto nella composizione dei componenti o nella gestione dello stato non
viene rilevato. &
Controllo dei tipi e analisi statica sull'intero client; verifica manuale dei
percorsi principali. Il client non contiene regole di dominio: la logica che
conta è verificata sul server. \\

Le prove di sistema sono manuali. &
Possono essere omesse sotto pressione, ed è il momento in cui servirebbero
di più. &
Sono elencate come casi di test con procedura e oracolo espliciti, non come
lista di controllo informale. \\

La verifica non comprende prove di carico o di sicurezza offensive. &
Il comportamento sotto carico e la resistenza ad attacchi non elementari non
sono noti. &
Fuori dagli obiettivi del progetto; le contromisure adottate sono
documentate nel SDD e verificate a livello funzionale. \\

\end{longtable}

\section{Interventi proposti}

Gli interventi che seguono non sono difetti ma direzioni di miglioramento
della verifica, ordinate per rapporto fra beneficio e costo.

\begin{enumerate}
  \item \textbf{Test di integrazione su un database effimero.} Eseguire le
    prove di sistema automaticamente su un'istanza MySQL creata e distrutta
    dalla catena di integrazione chiuderebbe il limite principale. È il
    singolo intervento con il beneficio maggiore.
  \item \textbf{Verifica automatica dei percorsi principali del client.} Una
    manciata di prove end-to-end su accesso, svolgimento di un quiz e
    inserimento di un feedback coprirebbe le regressioni più probabili
    nell'interfaccia.
  \item \textbf{Casi di test sui trigger.} Verificabili con asserzioni SQL
    dirette, una volta disponibile un database nella catena di integrazione.
  \item \textbf{Analisi di mutazione.} La copertura dice quali righe sono
    eseguite, non se le asserzioni sono in grado di rilevare un cambiamento
    di comportamento. Un'analisi di mutazione sui servizi darebbe una misura
    più severa dell'adeguatezza dei casi di test.
\end{enumerate}
